% This table is generated by analysis/analyze_manuscript_results.py.
% Do not edit numerical values manually.
\begin{table}[t!]
\centering
\caption{Separate comparison of the direct precedence edge set $E$ and its sparse augmentation $E^+$.}
\label{tab:edge-augmentation-ablation}
\footnotesize
\renewcommand{\arraystretch}{1.12}
\setlength{\tabcolsep}{5pt}
\begin{threeparttable}
\begin{tabular*}{\textwidth}{@{\extracolsep{\fill}} l c c c @{}}
\toprule
Threshold & Median $\Delta$ clauses (min--max) & Proof coverage $+ / = / -$ & Runtime $+ / = / -$ \\
\midrule
$Q^{\mathrm{top}}_{\max}$ & 0.58\%--1.09\% & 0 / 3 / 0 & 2 / 0 / 1 \\
$Q^{\mathrm{avg}}_{\max}$ & 0.58\%--1.14\% & 2 / 1 / 0 & 3 / 0 / 0 \\
\bottomrule
\end{tabular*}
\begin{tablenotes}[flushleft]
\scriptsize
\item Each row summarizes Prior CSE, SM, and SM-T in the separate $E^+$ test; it is not part of the Direct--Two-phase comparison in Table~\ref{tab:sat-main-results}. In the $+ / = / -$ columns, $+$ means that $E^+$ improves the metric, $=$ means no change, and $-$ means that it worsens the metric. Runtime uses the geometric mean on common proven cases; $\Delta$ clauses is computed pairwise from $E$ to $E^+$.
\end{tablenotes}
\end{threeparttable}
\end{table}
